\documentclass[12pt,a4paper]{article}
\usepackage{amsmath,amsfonts,amssymb}
\usepackage{graphicx}
\usepackage{hyperref}
\usepackage{enumitem}
\usepackage{geometry}
\geometry{margin=2.5cm}
\emergencystretch=2em

\begin{document}

\begin{center}
    {\Large \textbf{Master Thesis Project in Computational Physics}}\\[0.3cm]
    {\large \textbf{Parametric Matrix Models for the Black--Scholes Equation: Surrogate Models, Reduced-Order Operators and Options without Closed-Form Solutions}}\\[0.2cm]
    \vspace{0.3cm}
    {Duration: One Year (60 ECTS)}
\end{center}

\vspace{1cm}

\section*{1. Introduction}

The Black--Scholes partial differential equation (PDE)~[1,2] is one of
the central equations of quantitative finance.  It gives the fair
price of a European option as a function of time, of the price of the
underlying asset and of the volatility and the interest rate.  For
the simplest contracts the solution is known in closed form.  For
other contracts and for modified market assumptions the equation is
solved numerically, with finite differences or finite elements~[3].
In practice the price is needed again and again for different
parameters, for example when a model is calibrated to market data or
when the risk of a portfolio is computed, and this makes fast and
accurate surrogate models valuable.

Machine-learning surrogates for option prices have a long
history~[4].  Most of them are neural networks, which are not built to
respect the analytic properties of the PDE.  They need many
parameters and much data, their derivatives with respect to the
inputs (the sensitivities that are used for hedging) can be poor, and
they are unreliable outside the training region.

\emph{Parametric matrix models} (PMMs)~[5] are built from matrix
equations instead of neurons.  The trainable parameters are the
elements of matrices that depend on the input features, and the
outputs are eigenvalues of these matrices or expectation values in
their eigenvectors.  PMMs are universal function approximators, they
often need far fewer parameters than neural networks, their outputs
are analytic functions of the features, and they have extrapolated
well in the applications studied so far.  They are related to
eigenvector continuation and to reduced-basis methods~[6], in which a
large operator is projected on a small subspace, with the difference
that the small matrices are learned from data.

The Black--Scholes equation is well suited as a test of PMMs for
PDEs.  In suitable variables it is a Schr\"odinger equation in
imaginary time whose Hamiltonian depends linearly on the squared
volatility and on the interest rate~[7].  This is the structure a PMM
is designed to represent.  The closed-form solution gives exact
reference data, and the extensions of the model lead step by step to
problems where no closed form exists and a surrogate is of practical
use.

This thesis project develops and tests PMMs for the Black--Scholes
equation.  It has a common first part, with the PMM theory, a PMM
code and a finite-difference reference solver, and three possible
paths:
\begin{itemize}
  \item[\textbf{A.}] PMMs as surrogate models for the price and its
  sensitivities, benchmarked against neural networks and other
  regression methods;
  \item[\textbf{B.}] PMMs that follow the operator structure of the
  equation, as learned reduced-order models, compared with
  projection-based reduced-basis methods;
  \item[\textbf{C.}] PMMs for options without closed-form solutions:
  American options, several underlying assets, and local or stochastic
  volatility.
\end{itemize}
A thesis is expected to carry out the common part and Path~A or
Path~B in depth, and to apply the resulting model to at least one of
the problems of Path~C.  The choice is made together with the
supervisors after the first months.

\section*{2. The Black--Scholes Equation, Parametric Matrix Models and Scientific Motivation}

\textbf{The equation.}  The price $V(S,t)$ of a European option on an
asset with price $S$, volatility $\sigma$ and risk-free interest rate
$r$ obeys
\[
\frac{\partial V}{\partial t}
+ \frac{1}{2}\sigma^2 S^2 \frac{\partial^2 V}{\partial S^2}
+ r S \frac{\partial V}{\partial S}
- rV = 0,
\]
with the payoff as terminal condition, $V(S,T)=\max(S-K,0)$ for a
call option with strike price $K$ and maturity $T$.  The solution is
\[
  V=S\,N(d_1)-K e^{-r\tau}N(d_2),\qquad
  d_{1,2}=\frac{\ln(S/K)+(r\pm\tfrac12\sigma^2)\tau}{\sigma\sqrt{\tau}},
\]
where $\tau=T-t$ and $N$ is the cumulative normal distribution.  The
ratio $V/K$ depends on three dimensionless variables only,
$\ln(S/K)$, $\sigma^2\tau$ and $r\tau$.  The derivatives of the price
with respect to $S$, $\sigma$, $t$ and $r$, the so-called Greeks,
determine how an option is hedged and are as important as the price
itself.

\textbf{Operator form.}  With $x=\ln(S/K)$ and the time to maturity
$\tau$ the equation becomes
\[
  \frac{\partial V}{\partial\tau}=-H_{\rm BS}V,\qquad
  H_{\rm BS}=-\frac{\sigma^2}{2}\frac{\partial^2}{\partial x^2}
  +\Bigl(\frac{\sigma^2}{2}-r\Bigr)\frac{\partial}{\partial x}+r ,
\]
so that $V(\tau)=e^{-\tau H_{\rm BS}}V(0)$, with the payoff as initial
condition~[7].  Two properties of $H_{\rm BS}$ matter for what
follows.  It is affine in the parameters,
\[
  H_{\rm BS}=\sigma^2A+rB,\qquad
  A=-\frac12\frac{\partial^2}{\partial x^2}+\frac12\frac{\partial}{\partial x},
  \quad B=1-\frac{\partial}{\partial x},
\]
and it is not Hermitian, but similar to a Hermitian operator,
\[
  H_{\rm BS}=e^{\alpha x}\Bigl(-\frac{\sigma^2}{2}\frac{\partial^2}{\partial x^2}
  +\gamma\Bigr)e^{-\alpha x},\qquad
  \alpha=\frac12-\frac{r}{\sigma^2},\quad
  \gamma=r+\frac{(\sigma^2/2-r)^2}{2\sigma^2},
\]
which is the transformation to the heat equation.  Barrier options
add boundary conditions and give a discrete spectrum, as for a
particle in a box, and a volatility that depends on the asset price
adds a potential.

\textbf{Parametric matrix models.}  A PMM consists of primary
matrices, whose elements are analytic functions of the input features
$\boldsymbol c$, and of secondary matrices.  Its outputs are
eigenvalues of the primary matrices and bilinear forms of their
eigenvectors with the secondary matrices.  Two types of model are
studied.
\begin{itemize}
  \item \textbf{Regression PMM (Path~A).}  The features are
  $\boldsymbol c=(\ln(S/K),\sigma^2\tau,r\tau)$, the primary matrix is
  $M(\boldsymbol c)=M_0+\sum_lc_lM_l$ with trainable Hermitian
  $n\times n$ matrices, and the output $V/K$ is a bilinear form
  $\boldsymbol v^\dagger\Delta\,\boldsymbol v$ of its eigenvectors
  with a trainable secondary matrix $\Delta$.  The model knows nothing about the PDE
  and is trained on prices only.
  \item \textbf{Operator PMM (Path~B).}  The model has the form of the
  solution,
  \[
    V(x_i,\tau;\sigma,r)\simeq\boldsymbol s_i^{T}
    \exp\bigl[-\tau\,(\sigma^2M_1+rM_2)\bigr]\boldsymbol\phi_0 ,
  \]
  with trainable $n\times n$ matrices $M_1$ and $M_2$ in the place of
  $A$ and $B$, a trainable vector $\boldsymbol\phi_0$ for the payoff
  and trainable vectors $\boldsymbol s_i$ for the points $x_i$ at
  which the price is wanted.  This is a reduced-order model of the
  PDE.  A projection-based reduced-basis method~[6,8] has the same
  form, with all quantities computed from the discretized operator and
  from snapshots of the solution, while the PMM learns them from
  prices.  The similarity transformation above suggests a variant with
  Hermitian matrices.
\end{itemize}
Since the outputs of both models are analytic in the features, the
Greeks follow from analytic derivatives of the model.

The methods are compared through the accuracy of the price and of the
Greeks, the number of training points, the number of trainable
parameters, the cost of an evaluation, and the behavior outside the
training region.  The project is organized around four questions:
\begin{enumerate}
  \item How small can a PMM be, and how little data does it need, to
  reproduce the price and the Greeks to a given accuracy over the
  parameter domain, compared with neural networks, physics-informed
  neural networks, Gaussian processes and kernel methods?
  \item How well does it extrapolate to large and small volatility,
  to options deep in or out of the money and to long maturities, and
  does it respect the constraints that exclude arbitrage (bounds on
  the price, monotonicity in $S$, convexity in $K$)?
  \item Does a model that follows the operator structure need less
  data and extrapolate better than the regression PMM, and how does
  it compare with a reduced basis obtained by projection?
  \item Do the results carry over to problems without a closed-form
  solution, where a surrogate is needed in practice?
\end{enumerate}

\section*{3. Objectives of the Thesis}

\subsection*{Objective 1: PMMs, reduced-order methods and the reference solver (common part)}
\begin{itemize}
    \item Study and summarize the PMM framework of Ref.~[5], and its
    connections to eigenvector continuation, reduced-basis methods and
    implicit-function models~[6].
    \item Derive the Black--Scholes equation, its analytical solution
    and the Greeks, the transformation to the heat equation and the
    operator form of Section~2.
    \item Implement a stable finite-difference solver
    (Crank--Nicolson) as reference, test its convergence against the
    analytical solution, and generate data sets over volatilities,
    interest rates, strikes and maturities.
    \item Implement a PMM code with the analytic gradients of
    Ref.~[5] or with automatic differentiation, and test small models
    ($5\times5$, $7\times7$) on simplified payoff functions and on
    models with one parameter.
\end{itemize}

\subsection*{Objective 2 (Path A): PMM surrogates for prices and Greeks}
\begin{itemize}
    \item Construct PMMs with the input features $(S,t,\sigma,r,K)$
    and with the lower-dimensional encoding
    $(\ln(S/K),\sigma^2\tau,r\tau)$, and explore Hermitian and unitary
    formulations.
    \item Train the models on option prices and study the accuracy as
    a function of the matrix dimension and of the number and placement
    of the training points.
    \item Compute the Greeks from the trained models and compare with
    the analytical expressions.
    \item Compare with feed-forward neural networks, physics-informed
    neural networks~[9], Gaussian processes and kernel ridge
    regression~[10] in accuracy, computational efficiency, number of
    trainable parameters and quality of the extrapolation.
    \item Study the extrapolation to large and small volatility, deep
    in-the-money and out-of-the-money regimes and long maturities,
    and test the no-arbitrage constraints.
\end{itemize}

\subsection*{Objective 3 (Path B): PMMs with the operator structure of the equation}
\begin{itemize}
    \item Implement the operator PMM of Section~2, with the matrix
    exponential and its derivatives, in the non-Hermitian form and in
    the Hermitian form after the similarity transformation.
    \item Investigate how the payoff and the boundary conditions can
    be embedded in the vectors and secondary matrices of the model,
    for call and put options and for barrier options.
    \item Build a projection-based reduced-order model from the
    finite-difference operator and snapshots of the solution~[8], and
    compare the matrices it gives with the matrices the PMM has
    learned, and their spectra with the spectrum of the discretized
    operator.
    \item Compare the operator PMM with the regression PMM of Path~A
    and with the projection-based model in data requirements and in
    the extrapolation in maturity, volatility and interest rate.
    \item Study the region close to maturity and at the money, where
    the kink of the payoff limits any smooth model.
\end{itemize}

\subsection*{Objective 4 (Path C): Options without closed-form solutions}
\begin{itemize}
    \item American options: add the early-exercise condition to the
    finite-difference solver with a penalty method~[11], and extend
    the PMMs to the resulting nonlinear problem.
    \item Options on two assets: a two-dimensional PDE with the
    correlation as an additional parameter.
    \item Local and stochastic volatility~[12]: a volatility that
    depends on the asset price, and the Heston model with its
    two-dimensional PDE and five parameters.
    \item Use a trained PMM as a fast pricing function in the
    calibration of model parameters to option prices, and compare
    with neural-network surrogates~[4].
\end{itemize}

\subsection*{Objective 5: Performance analysis and interpretation (all paths)}
\begin{itemize}
    \item Quantify accuracy, data requirements, number of parameters
    and evaluation cost of the PMMs against the reference solver and
    the baselines.
    \item Analyze the learned matrices: their spectra, the analytic
    structure of the outputs in the complex plane of the features,
    and how these limit the extrapolation.
    \item Establish where the models fail and how the failures can be
    detected without reference data.
\end{itemize}

\subsection*{Objective 6: Possible extensions}
\begin{itemize}
    \item The inverse problem: PMMs for the implied volatility as a
    function of the price.
    \item Other PDEs of the same type in physics and engineering
    (diffusion and Fokker--Planck equations), with the methods of this
    thesis.
    \item Tests on market data for index options.
\end{itemize}

\section*{4. Methods and Tools}

\begin{itemize}
  \item \textbf{PDE solvers:} Python with NumPy and SciPy for the
  finite-difference solvers, with C++ for the two-dimensional
  problems if needed.
  \item \textbf{PMMs and baselines:} PyTorch or JAX for the PMMs, with
  automatic differentiation through eigenvalue problems and matrix
  exponentials; Scikit-Learn and PyTorch for the baselines.
  \item \textbf{Useful background:} computational physics (FYS3150 or
  FYS4411), machine learning (FYS-STK4155), and basic knowledge of
  quantum mechanics and linear algebra.  No previous knowledge of
  finance is required.  The lecture notes of Ref.~[13] cover the
  numerical solution of PDEs and the machine-learning methods.
\end{itemize}
All codes and results will be documented in Jupyter notebooks so that
the thesis can be reproduced.

\section*{5. Work Plan and Milestones (One-Year Timeline)}

\subsection*{Semester 1 (Months 1--6)}

\textbf{Months 1--2: Literature study}
\begin{itemize}
    \item Read the PMM paper and supporting literature on
    reduced-basis methods.
    \item Study the derivation and the numerical properties of the
    Black--Scholes equation.
    \item Begin preliminary numerical experiments.
\end{itemize}
\textbf{Deliverable:} A short report summarizing PMMs and the
Black--Scholes equation.

\bigskip

\textbf{Months 3--4: Numerical solvers and data generation}
\begin{itemize}
    \item Implement the finite-difference solver and test grid
    convergence and accuracy.
    \item Generate data sets for a range of volatilities, interest
    rates and maturities.
    \item Implement the PMM code and test small models.
\end{itemize}
\textbf{Deliverable:} Verified PDE solver, data sets and a first
working PMM for models with one parameter.

\bigskip

\textbf{Months 5--6: PMM construction}
\begin{itemize}
    \item Path~A: regression PMMs over the full parameter domain;
    first comparison with a neural network.
    \item Path~B: operator PMM in non-Hermitian and Hermitian form;
    projection-based reduced-order model.
\end{itemize}
\textbf{Deliverable:} Working PMM for European options along the
chosen path.

\subsection*{Semester 2 (Months 7--12)}

\textbf{Months 7--9: Full training and analysis}
\begin{itemize}
    \item Assess accuracy, convergence and extrapolation; optimize
    matrix sizes and training.
    \item Path~A: Greeks and no-arbitrage constraints.  Path~B:
    barrier options and comparison of the learned and the projected
    matrices.
    \item Start of one problem of Path~C.
\end{itemize}
\textbf{Deliverable:} Fully trained PMMs for Black--Scholes pricing.

\bigskip

\textbf{Months 10--11: Benchmarking and applications}
\begin{itemize}
    \item Comparison with neural networks and other regression
    models; computational speed and parameter efficiency.
    \item Results for the chosen problem of Path~C; one of the
    extensions of Objective~6, depending on progress.
\end{itemize}
\textbf{Deliverable:} Benchmarking report with figures and tables.

\bigskip

\textbf{Month 12: Writing and final submission}
\begin{itemize}
    \item Completion of the thesis manuscript (methods and theory
    chapters are written continuously during the year) and
    preparation of the final presentation.
\end{itemize}
\textbf{Deliverable:} Completed thesis and defense presentation.

\section*{6. Expected Outcomes}

\begin{itemize}
    \item A documented implementation of PMMs for the Black--Scholes
    equation, with the reference solvers and data sets.
    \item A demonstration of what PMMs can do as surrogate PDE
    solvers, for prices and for sensitivities.
    \item Quantitative comparisons with neural networks, other
    regression methods, reduced-basis methods and numerical PDE
    solvers.
    \item An assessment of whether building the operator structure of
    a PDE into the model improves accuracy and extrapolation.
    \item First results for an option-pricing problem without a
    closed-form solution.
    \item A master thesis and, depending on the results, a
    contribution to a scientific article.
\end{itemize}

\section*{7. References}

\begin{enumerate}
    \item F.~Black and M.~Scholes, \textit{The Pricing of Options and Corporate Liabilities}, J.\ Polit.\ Econ.\ \textbf{81}, 637 (1973); R.~C.~Merton, \textit{Theory of Rational Option Pricing}, Bell J.\ Econ.\ Manag.\ Sci.\ \textbf{4}, 141 (1973).
    \item J.~C.~Hull, \textit{Options, Futures, and Other Derivatives}, Pearson; P.~Wilmott, S.~Howison and J.~Dewynne, \textit{The Mathematics of Financial Derivatives}, Cambridge University Press (1995).
    \item D.~Tavella and C.~Randall, \textit{Pricing Financial Instruments: The Finite Difference Method}, Wiley (2000).
    \item J.~M.~Hutchinson, A.~W.~Lo and T.~Poggio, \textit{A Nonparametric Approach to Pricing and Hedging Derivative Securities Via Learning Networks}, J.\ Finance \textbf{49}, 851 (1994); J.~Ruf and W.~Wang, \textit{Neural networks for option pricing and hedging: a literature review}, J.\ Comput.\ Finance \textbf{24}, 1 (2020); B.~Horvath, A.~Muguruza and M.~Tomas, \textit{Deep learning volatility: a deep neural network perspective on pricing and calibration in (rough) volatility models}, Quant.\ Finance \textbf{21}, 11 (2021).
    \item P.~Cook, D.~Jammooa, M.~Hjorth-Jensen, D.~D.~Lee and D.~Lee, \textit{Parametric matrix models}, Nat.\ Commun.\ \textbf{16}, 5929 (2025).
    \item D.~Frame, R.~He, I.~Ipsen, D.~Lee, D.~Lee and E.~Rrapaj, \textit{Eigenvector continuation with subspace learning}, Phys.\ Rev.\ Lett.\ \textbf{121}, 032501 (2018); T.~Duguet, A.~Ekstr\"om, R.~J.~Furnstahl, S.~K\"onig and D.~Lee, \textit{Colloquium: Eigenvector continuation and projection-based emulators}, Rev.\ Mod.\ Phys.\ \textbf{96}, 031002 (2024); A.~Quarteroni, A.~Manzoni and F.~Negri, \textit{Reduced Basis Methods for Partial Differential Equations}, Springer (2016).
    \item B.~E.~Baaquie, \textit{Quantum Finance: Path Integrals and Hamiltonians for Options and Interest Rates}, Cambridge University Press (2004).
    \item R.~Cont, N.~Lantos and O.~Pironneau, \textit{A Reduced Basis for Option Pricing}, SIAM J.\ Financial Math.\ \textbf{2}, 287 (2011).
    \item M.~Raissi, P.~Perdikaris and G.~E.~Karniadakis, \textit{Physics-informed neural networks: A deep learning framework for solving forward and inverse problems involving nonlinear partial differential equations}, J.\ Comput.\ Phys.\ \textbf{378}, 686 (2019).
    \item C.~E.~Rasmussen and C.~K.~I.~Williams, \textit{Gaussian Processes for Machine Learning}, MIT Press (2006).
    \item P.~A.~Forsyth and K.~R.~Vetzal, \textit{Quadratic Convergence for Valuing American Options Using a Penalty Method}, SIAM J.\ Sci.\ Comput.\ \textbf{23}, 2095 (2002).
    \item B.~Dupire, \textit{Pricing with a smile}, Risk \textbf{7}, 18 (1994); S.~L.~Heston, \textit{A Closed-Form Solution for Options with Stochastic Volatility with Applications to Bond and Currency Options}, Rev.\ Financ.\ Stud.\ \textbf{6}, 327 (1993).
    \item M.~Hjorth-Jensen, \textit{Computational Physics} and \textit{Applied Data Analysis and Machine Learning}, lecture notes and codes,\\ \url{https://github.com/CompPhysics/ComputationalPhysics} and\\ \url{https://github.com/CompPhysics/MachineLearning}.
\end{enumerate}

\end{document}
